\section{Introduction}
\label{sec:intro}

Builders of LLM-agent systems increasingly design their products around MCP
servers: file stores, databases, web services, developer platforms, and other
external tools exposed behind a uniform calling interface. For such a system,
the practical question is not whether an agent can issue a tool call in
isolation, but whether the configured agent--server stack accomplishes the
user's goal across many real, interacting services. Yet the benchmarks used to
answer that question mostly score one of two proxies: the agent's final
natural-language answer, matched against a reference string, or the set of
tools the agent was expected to call. \textbf{Both are fragile precisely where
deployment matters.} A reference answer goes stale the moment it depends on
live data---today's price, this morning's commit, the current contents of a
table---and a fixed tool list cannot be recovered from the prompt alone, so it
accumulates spurious ``required'' tools that penalize an agent for taking a
different, equally valid route.

The companion methodology paper, DynamicMCPBench
\citep{dynamicmcpbench2026}, argues that the organizing object of the
benchmark should be neither the answer nor the tool list but the
\emph{execution trace}: explore a goal live, distill the successful
trajectory into path-agnostic \emph{effect checkpoints}, and grade a
candidate purely on whether it reproduces those effects. That argument is
easy to state and hard to feel. A reader of the paper sees tables; what they
do not see is the moment a confident, fluent answer is marked
\textsf{FAILED} because the agent never produced a required piece of
evidence, or the moment a ``wrong'' answer is marked \textsf{SOLVED} because
the live data simply moved.

\begin{figure*}[tbp]
  \centering
  % Main system figure (full-width; \input from a figure* in introduction.tex).
% TODO(figure): replace \placeholderfig with an annotated screenshot of the
% scoring stage, e.g. \includegraphics[width=\textwidth]{figures/fig_studio.png}
% captured from the running studio (dmcp_studio.html / deployed instance),
% with callouts on the Effect/Answer toggle, the verdict bar, and an editable
% equivalence-set chip.
\placeholderfig[\textwidth]{studio screenshot --- scoring stage with the
  Effect/Answer toggle and the flipped verdict bar (TODO: drop in PNG)}

  \caption{\textbf{DMCP Studio} (scoring stage, REPLAY mode). The four-stage
    stepper (collect $\rightarrow$ explore $\rightarrow$ distill $\rightarrow$
    score) and the LIVE/REPLAY toggle sit in the header. Here the
    \textsf{hermes3-8b} candidate is marked \textsf{FAILED}: its trajectory
    (left) reproduces every effect \emph{except} the income-statement
    checkpoint (\#5, red in the ledger, right), even though its prose summary
    would pass answer-matching---the incomplete-aggregation case called out in
    the note below. The Effect\,$\rightleftharpoons$\,Answer toggle re-grades
    the \emph{same} recorded run two ways without a second backend call, and the
    \texttt{download}\,/\,\texttt{get\_price\_history} equivalence set on
    checkpoint~\#3 is an editable chip the visitor can toggle to re-score live.}
  \label{fig:studio}
\end{figure*}

\textbf{DMCP Studio} makes benchmark design and execution interactive. It wraps
the existing DynamicMCPBench pipeline, helps the user decide what benchmark is
statistically worth generating, and then walks that benchmark through four live
stages---\emph{collect} servers, \emph{explore} a generated goal into a
reference trace, \emph{distill} that trace into a TaskSpec of effect
checkpoints, and \emph{score} a candidate against them
(Figure~\ref{fig:studio}). The pre-run \emph{design} stage is not a passive
explanation panel: Benchmark Advisor computes what claim the planned corpus can
support, what task budget it requires, and what assumptions must remain visible
before any corpus is generated (Stage~0, \S\ref{sec:advisor}). The interactive
centerpiece for scoring remains the
\textbf{Effect\,$\rightleftharpoons$\,Answer toggle}, which grades one
candidate run under both regimes and flips the verdict on screen. Three curated
cases make the disagreement concrete: (1)~a clean pass that reaches every effect
through a \emph{different but equivalent} tool; (2)~a confident agent that
passes answer-matching but \emph{fails} effect-scoring because it skipped
required work (incomplete aggregation); and (3)~a correct agent that
\emph{fails} answer-matching but passes effect-scoring because the live data
drifted since the reference was recorded.

Two design choices keep the demo both honest and useful as a system. First, a
\textbf{LIVE/REPLAY} mode switch: REPLAY serves frozen fixtures
deterministically---this is at once a booth-safety measure and on-thesis, since
the methodology itself scores under deterministic replay---while LIVE drives the
real pipeline against read-only servers as proof. Second, the studio never
reimplements scoring: all real work flows through thin adapters over the
published pipeline and the Advisor services, so what the visitor sees is a
usable system rather than a mock-up.

\paragraph{Audience.} DMCP Studio targets three groups: agent builders who
need a deployment-specific diagnostic on their own MCP fleet, evaluation
researchers comparing scoring methodologies, and practitioners who want to
understand \emph{why} answer-matching misleads on live, stateful tools before
adopting an effect-based alternative.

\paragraph{Contributions.} (i)~An interactive benchmark-design and execution
studio that exposes a trace-grounded, answer-agnostic evaluation pipeline
end-to-end in the browser, with a verdict-flipping Effect/Answer toggle as its
central scoring interaction (\S\ref{sec:system}); (ii)~Benchmark Advisor, a
pre-generation statistical design assistant that maps user intent to a scoped
benchmark plan, searches candidate budgets and effect targets, pre-computes the
MDE implied by the unique-task budget, and carries only defensible plans toward
corpus generation (\S\ref{sec:advisor}); and (iii)~a LIVE/REPLAY design plus a
system evaluation:
studio-vs-batch verdict agreement, per-stage latency, and an informal usability
pass (\S\ref{sec:eval}).
